% !TeX program = pdflatex
%==============================================================================
%  Lista Teórica 06 --- Pré-processamento e Pipelines
%  O gabarito sai de "Lista teorica 06 - gabarito.tex".
%==============================================================================
\providecommand{\gabaritoopt}{}
\documentclass[11pt,a4paper]{article}
\usepackage[\gabaritoopt]{../../recursos/latex/estilo-lista}

\begin{document}
\cabecalholista{06}{Pré-processamento e \emph{Pipelines}}

%------------------------------------------------------------------------------
\begin{exercicio}[revisando o notebook de um colega]
Um colega lhe manda o notebook dele na véspera de entregar o relatório e pede uma
olhada. Ele leu a Aula~06 na diagonal, e o resultado é um pouco de cada coisa: em
alguns trechos fez exatamente o certo; em outros cometeu um deslize que não muda
nada; e em pelo menos um está prestes a reportar um número que não existe.

Para separar os três casos você tem uma pergunta só, e ela não é ``quanto essa
etapa aprendeu dos dados''. É: \textbf{essa etapa olhou o $Y$?}

Diga, para cada trecho, se o vazamento é \textbf{grave}, \textbf{leve}, ou se não
há vazamento nenhum --- e justifique em uma linha, aplicando a pergunta.
\begin{enumerate}[label=(\alph*),leftmargin=1.1cm]
  \item Ele padronizou as covariáveis com a média e o desvio do banco inteiro, e
        só depois separou treino e teste.
  \item De $10\,000$ covariáveis, ficou com as $50$ mais correlacionadas com $Y$
        --- correlação calculada no banco inteiro --- e rodou a validação cruzada
        só nessas $50$.
  \item Ajustou um PCA no banco inteiro, guardou as $20$ primeiras componentes, e
        aí separou treino e teste.
  \item Preencheu os faltantes com a mediana da própria coluna, calculada no banco
        inteiro.
  \item O banco tem uma linha por consulta, e o mesmo paciente aparece em várias
        delas. Ele rodou \texttt{KFold} sobre as linhas.
  \item Pôs o \texttt{StandardScaler} dentro de um \texttt{Pipeline}, e passou o
        \emph{pipeline} inteiro para o \texttt{GridSearchCV}.
\end{enumerate}
\end{exercicio}

\begin{solucao}
Vale agrupar por veredito, porque dentro de cada grupo o argumento é o mesmo:
\textbf{(a), (c) e (d) são leves}; \textbf{(b) e (e) são graves}; e \textbf{(f)
não é vazamento} --- é justamente a conduta correta.

\medskip
\textbf{Os três leves.} Todos mexem só em $\bm X$. O \texttt{StandardScaler} de
(a) calcula duas estatísticas por coluna, a média e o desvio; a imputação de (d)
calcula uma, a mediana; e o PCA de (c) procura direções de máxima variância. Em
nenhum dos três o $Y$ entra na conta, e é isso que os salva: sem olhar a resposta,
não existe caminho por onde fabricar acerto. Trocar o treino pelo banco inteiro
desloca essas estatísticas por algo da ordem de $1/\sqrt n$, e desloca igual para
todas as observações, independentemente do $y$ de cada uma.

O (c) é o que mais engana, e vale parar nele. Ajustar o PCA no banco todo
\emph{de fato} usa informação do teste --- as componentes saem outras. Mas ``usar
informação'' e ``inflar o número que se vai reportar'' são coisas diferentes, e
sem olhar o $Y$ só a primeira acontece. A Aula~E2 mediu quatro configurações com
$y$ de ruído puro, e em nenhuma delas o $R^2$ subiu.

\medskip
\textbf{Os dois graves, e por motivos distintos.} No (b), a escolha das $50$
colunas foi feita \emph{pela correlação com $Y$} --- inclusive com a parte do $Y$
que está nas dobras de validação. Quando a validação cruzada roda depois, cada
dobra encontra covariáveis que ela própria ajudou a escolher, e reencontra ali o
seu próprio ruído. Com $10\,000$ colunas disponíveis, o acaso tem espaço de sobra:
a Aula~06 mede $R^2 = +0{,}40$ num caso em que a resposta é ruído puro e o valor
verdadeiro é zero.

Já o (e) não é sobre pré-processamento nenhum. É a \emph{mesma unidade} dos dois
lados da divisão: consultas do mesmo paciente, umas no treino e outras na
validação. O modelo não precisa aprender nada sobre medicina --- basta reconhecer
o paciente, que ele já viu.

\medskip
\textbf{E o (f)}, que é o que se deve fazer: dentro do \emph{pipeline}, a
padronização é reajustada em cada dobra, com a média e o desvio de quem está no
treino \emph{daquela} dobra. Nada vaza, e não custou esforço nenhum.

\medskip
O par que vale guardar é este: \textbf{o \texttt{Pipeline} resolve (a), (b), (c) e
(d)} --- basta pôr cada etapa dentro dele, a seleção de variáveis inclusive ---,
\textbf{mas não faz nada pelo (e)}. Ele controla \emph{quando} cada etapa é
ajustada, não \emph{como} as linhas são repartidas, e não tem como saber que duas
delas são o mesmo paciente. Para o (e) a ferramenta é o \texttt{GroupKFold}, com o
identificador do paciente em \texttt{groups}. São duas perguntas independentes, e
você precisa das duas na cabeça: \emph{a etapa olha o $Y$?} e \emph{a divisão
respeita a unidade?}
\end{solucao}

%------------------------------------------------------------------------------
\begin{exercicio}[a receita inteira, num objeto só]
Você chega a um banco e recebe a tarefa de sempre: prever inadimplência. O arquivo
tem seis colunas. Três são numéricas --- \texttt{idade}, \texttt{renda} e
\texttt{saldo} ---, e a \texttt{renda} vem com buracos, porque nem todo cliente
informa. Duas são texto, \texttt{estado} e \texttt{profissao}, e essas vieram
completas. A sexta é a resposta, \texttt{inadimplente}, que vale $0$ ou $1$.

O que você quer fazer com elas é o que a aula descreveu: preencher os faltantes
das numéricas pela mediana, padronizá-las, transformar as duas categóricas em
indicadoras, e ajustar uma regressão logística com penalização $\ell_2$,
escolhendo o $C$ por validação cruzada de cinco dobras. E quer tudo isso sem que o
número que você vai apresentar no fim saia otimista.
\begin{enumerate}[label=(\alph*),leftmargin=1.1cm]
  \item Escreva o código, com \texttt{ColumnTransformer}, \texttt{Pipeline} e
        \texttt{GridSearchCV}.
  \item Se a grade de $C$ tem seis valores, quantas vezes a mediana da
        \texttt{renda} acaba sendo calculada ao longo de toda a busca? Justifique,
        e compare com o que aconteceria se você a imputasse uma vez só, no começo.
  \item Por que a ordem tem de ser imputar \emph{e depois} padronizar, e não o
        contrário?
\end{enumerate}
\end{exercicio}

\begin{solucao}
\textbf{(a)} As colunas pedem tratamentos diferentes, então o
\texttt{ColumnTransformer} abre dois ramos --- e o ramo numérico é, ele próprio,
um \emph{pipeline} de duas etapas:
\begin{lstlisting}
import numpy as np
import sklearn.model_selection as skm
from sklearn.compose import ColumnTransformer
from sklearn.impute import SimpleImputer
from sklearn.linear_model import LogisticRegression
from sklearn.pipeline import Pipeline
from sklearn.preprocessing import OneHotEncoder, StandardScaler

numericas   = ["idade", "renda", "saldo"]
categoricas = ["estado", "profissao"]

prep = ColumnTransformer([
    ("num", Pipeline([("imp",   SimpleImputer(strategy="median")),
                      ("escala", StandardScaler())]), numericas),
    ("cat", OneHotEncoder(handle_unknown="ignore"), categoricas),
])

tubo = Pipeline([("prep", prep),
                 ("modelo", LogisticRegression(max_iter=2000))])

busca = skm.GridSearchCV(
    tubo,
    {"modelo__C": np.logspace(-3, 2, 6)},
    cv=skm.StratifiedKFold(5, shuffle=True, random_state=0),
    scoring="roc_auc",
).fit(X, y)
\end{lstlisting}

Três escolhas do código merecem uma palavra. O \texttt{handle\_unknown="ignore"}
evita que o ajuste quebre quando uma profissão rara aparece só na validação e o
codificador a vê pela primeira vez --- é o tipo de erro que não acontece no seu
notebook e acontece em produção, seis meses depois. O \texttt{StratifiedKFold}
mantém a proporção de inadimplentes em cada dobra: em classificação ele já é o
padrão do \texttt{GridSearchCV}, mas escrevê-lo deixa o \texttt{shuffle} sob o seu
controle, e não sob o de quem ordenou o arquivo (Exercício~3 da Lista
prática~03). E a penalização $\ell_2$ que o enunciado pede \textbf{não precisa ser
escrita}: é o padrão da \texttt{LogisticRegression}. Escrever
\texttt{penalty="l2"} não está errado, mas o argumento foi depreciado na versão
1.8 do \texttt{scikit-learn} e desaparece na 1.10 --- omiti-lo funciona em
qualquer versão.

\smallskip
\textbf{(b)} \textbf{Trinta vezes}: seis valores de $C$ vezes cinco dobras. Para
cada combinação, o \texttt{GridSearchCV} chama o \texttt{fit} do \emph{pipeline
inteiro} usando só os $4/5$ de treino daquela dobra, e ali dentro o
\texttt{SimpleImputer} recalcula a mediana do zero. (Com \texttt{refit=True}, que
é o padrão, ainda há um $31^{\text{o}}$ ajuste no conjunto completo ao final, para
construir o \texttt{best\_estimator\_}.)

Comparado com imputar uma vez só no começo, isso é trinta vezes mais trabalho para
chegar à mesma coluna preenchida --- ou quase a mesma, e é justamente o
\emph{quase} que interessa. A mediana muda de dobra para dobra, e essa variação faz
parte do que você está tentando medir: o desempenho não é do modelo sozinho, é do
\emph{procedimento} que imputa, padroniza e ajusta. Imputar antes tira essa
variação da conta e usa, em cada dobra, uma mediana que as observações de
validação ajudaram a calcular.

Pelo critério do Exercício~1 isso é um vazamento \emph{leve} --- a mediana não olha
o $Y$ ---, e o efeito numérico seria pequeno. A questão é que corrigi-lo custou
zero: já estava escrito.

\smallskip
\textbf{(c)} Pelo motivo prosaico primeiro: se você padronizar antes, o
\texttt{StandardScaler} encontra \texttt{NaN} na coluna \texttt{renda} e a média e
o desvio saem \texttt{NaN} também, contaminando a coluna inteira. Na melhor das
hipóteses o código quebra na sua frente; na pior ele segue em silêncio com uma
coluna de \texttt{NaN}.

E há o motivo conceitual, que sobrevive mesmo se a implementação fosse esperta o
bastante para ignorar os faltantes. Imputar depois de padronizar significa inserir
a mediana \emph{da escala já padronizada}: você acabou de deixar a coluna com média
$0$ e desvio $1$, e aí acrescenta um punhado de valores que desloca as duas coisas.
A padronização que você fez deixa de valer para os dados que efetivamente chegam ao
modelo. Imputando primeiro, a coluna chega completa ao \texttt{StandardScaler}, e a
média e o desvio que ele calcula são os dos números que o modelo vai de fato ver.
\end{solucao}


\end{document}
